% TOP_PROBLEM: {"id": 3100064, "title": "Show that there is some B so that no integer appears more than B times among the binomial coefficients", "category_id": 1, "category": {"id": 1, "name": "number_theory", "display_name": "Number Theory", "description": "Properties of integers, prime numbers, Diophantine equations.", "slug": "number-theory", "order_index": 1, "created_at": "2026-07-31T15:26:25.670Z"}, "status": "partially_solved"}
% FIRST_POSTED: null
% ENTRY_KIND: statement_only
% Imported from: batch11.tex; UnsolvedMath ID 3100064
\documentclass[11pt]{article}
\usepackage{amsmath,amssymb,mathtools}
\usepackage{fontspec,unicode-math}
\setmainfont{Latin Modern Roman}
\setmathfont{Latin Modern Math}
\usepackage{xurl,hyperref}
\newcommand{\sref}[2]{\hyperlink{source:#1:#2}{[S#2]}}
\newcommand{\cataloguescope}[1]{\paragraph{Mathematical statement.}#1}
\begin{document}
% UnsolvedMath ID 3100064; source catalogue code AMR-030-0064
\setcounter{section}{3100063}
\section{Show that there is some B so that no integer appears more than B times among the binomial coefficients}\label{3100064}
{\small\sffamily UnsolvedMath ID 3100064 · Number Theory}
\subsection{Definitions and mathematical statement}

\paragraph{Shared notation.}$\mathbb N=\{1,2,\ldots\}$, and $\mathbb Z,\mathbb Q,\mathbb R,\mathbb C$ denote the integers, rationals, reals and complex numbers. Any other conventions are specified locally.


For integers $n\ge0$ and $0\le m\le n$, define $\binom nm=n!/(m!(n-m)!)$. For an integer $t\ge2$, let
\[N(t)=\#\{(n,m)\in\mathbb Z^2:1\le m<n,\ \binom nm=t\}.\]
The symmetric positions $m$ and $n-m$ are counted separately. The boundary entries $\binom n0=\binom nn=1$ occur infinitely often; Singmaster's conjecture concerns values greater than one, as specified in the original statement cited below.

\paragraph{Statement recorded in the supplied source.}
Is there an absolute finite constant $B$ such that $N(t)\le B$ for every integer $t\ge2$? \sref{3100064}{2}

\subsection{Short English statement}
Is there a uniform bound on the number of occurrences in Pascal’s triangle of any integer greater than one?
\subsection{Sources}
\begin{description}
\item[\hypertarget{source:3100064:1}{S1}] ULAM.AI and the UnsolvedMath Contributors, \emph{UnsolvedMath: A Curated Collection of Open Mathematics Problems}, record 3100064 (AMR-030-0064). CC BY 4.0. \url{https://huggingface.co/datasets/ulamai/UnsolvedMath}.
\item[\hypertarget{source:3100064:2}{S2}] Kaisa Matomäki, Maksym Radziwiłł, Xuancheng Shao, Terence Tao and Joni Teräväinen, \emph{Singmaster\textquotesingle s conjecture in the interior of Pascal\textquotesingle s triangle} (2021), Section1, Conjecture1.1 and the following boundary-entry convention, page1. \url{https://arxiv.org/abs/2106.03335}.
\end{description}
\label{end:3100064}
\end{document}
